\section{Unexpected and Impossible Exceptions}\label{sec:ImpossibleExceptions}
The interface \lstinline|java.lang.Appendable|\index{java.lang!Appendable} \autocite{ORACLE_DOC_APPENDABLE_INTERFACE} is implemented by several classes, including
\begin{itemize}[nosep]
	\item{\lstinline|java.lang.StringBuilder|\index{java.lang!StringBuilder} \autocite{ORACLE_DOC_STRINGBUILDER_CLASS}}
	\item{\lstinline|java.io.StringWriter|\index{java.io!StringWriter} \autocite{ORACLE_DOC_STRINGWRITER_CLASS}}
	\item{\lstinline|java.io.PrintWriter|\index{java.io!PrintWriter} \autocite{ORACLE_DOC_PRINTWRITER_CLASS}}
	\item{\lstinline|java.io.PrintStream|\index{java.io!PrintStream} \autocite{ORACLE_DOC_PRINTSTREAM_CLASS}}
\end{itemize} 

The documentation for \lstinline|Appendable::append|\index{java.lang!Appendable!append()} looks like this \autocite{ORACLE_DOC_APPENDABLE:append}:
\begin{center}
	\includegraphics[scale=0.55]{coding/errorHandling/Appendable_append}
\end{center}

You see that the method \lstinline|appendable()|\index{java.lang!Appendable!append()} declares to throw an \lstinline|IOException|\index{java.io!IOException} in case of an error.

This exception in not declared for the implementation of \lstinline|appendable()|\index{java.lang!Appendable!append()} by neither \lstinline|java.lang.StringBuilder|\index{java.lang!StringBuilder} nor \lstinline|java.io.PrintStream|\index{java.io!PrintStream}.

\keeplisting{3}
Consequently, the code snippets
\begin{lstlisting}
public final void toStringBuilder()
{
    final StringBuilder stringBuilder = new StringBuilder();
    stringBuilder.append( "Text" );
}   //  toStringBuilder()
\end{lstlisting}
and
\keeplisting{3}
\begin{lstlisting}
public final void toPrintStream()
{
    final PrintStream printStream = new PrintStream( out );
    printStream.append( "Text" );
}   //  toPrintStream()
\end{lstlisting}
will compile without a problem.

\keeplisting{5}
And now this code snippet:
\begin{lstlisting}
final StringBuilder stringBuilder = new StringBuilder();
final PrintStream printStream = new PrintStream( out );

final Appendable appendable = toConsole ? printStream : stringBuilder;
try
{
    appendable.append( "Text" );
}
catch( final IOException e )
{
    // What's up here?
}
\end{lstlisting}

Despite neither \lstinline|PrintStream::append|\index{java.io!PrintStream!append()} nor \lstinline|StringBuilder::append|\index{java.lang!StringBuilder!append()} may throw an exception\index{exception}, the Java syntax requires the \lstinline|catch|\index{throw!catch} block starting in line~50.

What should we do in this \lstinline|catch|\index{throw!catch} block?

Of course, we can place a comment here and leave it otherwise empty. Or we just print or log the \lstinline|IOException|\index{java.io!IOException}.

\keeplisting{4}
My recommendation is to throw an error in line~11, either
\begin{lstlisting}[firstnumber=9]
catch( final IOException e )
{
    throw new AssertionError( "Impossible Exception encountered", e );
}
\end{lstlisting}

or you create an \lstinline|ImpossibleExceptionError| for this purpose and throw that. Refer to chapter \tqfullvref{sec:UsefulCustomExceptions}. This approach protects you from API changes in later versions of the \Java Runtime Library, as the related error will not remain unnoticed.

\keeplisting{10}
Another example from this category is this one:
\begin{lstlisting}
final String s;
try
{
    s = new String( byteArray, "UTF-8" );
}
catch( final UnsupportedEncodingException e )
{
    throw new ImpossibleExceptionError( "UTF-8 must exist!", e );
}
\end{lstlisting}

The documentation for \lstinline|java.nio.charset.Charset|\index{java.nio.charset!Charset} contains a list of standard charsets \autocite{ORACLE_DOC_CHARSET:StandardCharsets} that must be supported by each JDK/JVM implementation. According to that list, the UTF-8\index{UTF-8} \autocite{WIKIPEDIA:UTF8} encoding is mandatory and will be availabe always, and consequently, the code sequence above will \textit{never} throw an \lstinline|UnsupportedEncodingException|. 

Nevertheless, exceptions\index{exception} should never be swallowed. Instead we throw an error in the \lstinline|catch|\index{throw!catch} block, indicating the error condition. In the improbable case that UTF-8\index{UTF-8} is no longer supported by an upcoming \Java runtime environment, the problem will surface, and the program will not die silently without any hint about the reason. It also protects you from malicious changes to the name of the character encoding: a none existing \lstinline|Charset|\index{java.nio.charset!Charset} will cause noise.

Of course, the problem could be solved differently, by using an instance of \lstinline|java.nio.charset.Charset|\index{java.nio.charset!Charset} \autocite{ORACLE_DOC_CHARSET_CLASS} as the argument to the \lstinline|String|\index{java.lang!String} constructor, instead of the \textit{name} of the character encoding. For this case, the constructor does not declare an \lstinline|UnsupportedEncodingException|\index{java.nio.charset!UnsupportedEncodingException}.
 
\keeplisting{1}
That would look like this:
\begin{lstlisting}
final String s = new String( byteArray, StandardCharsets.UTF_8 );
\end{lstlisting}

%------------------------------------------------------------------------------

\chapter*{Road Work ahead!}


















Some APIs in the JDK declares a checked exception, but that will never be thrown for the given arguments by any means. Most famous for this are the conversions from byte arrays to Strings and back:



























It is also possible that you checked the preconditions, but you still have to deal with some checked exceptions:
\begin{lstlisting}
…
final var file = new File( … );
if( file.exists() )
{
    try( final var input = new FileInputStream( file ) )
    {
        …
    }
    catch( final FileNotFoundException e )
    {
        throw new UnexpectedExceptionError( "File was there when checked!", e );
    }
}
…
\end{lstlisting}

A sample for the class \lstinline|UnexpectedExceptionError| can be found in \autocite{TQUADRAT_ORG_FOUNDATION_UNEXPECTEDEXCEPTIONERROR}.

Another example is the interface \lstinline|java.lang.Appendable|\autocite{ORACLE_DOC_APPENDABLE_INTERFACE}: it is implemented by \lstinline|java.io.Writer| and \lstinline|java.io.PrintStream|, but also by \lstinline|java.lang.StringBuilder| (and a few more). The method \lstinline|java.lang.Appendable::append| declares to throw an \lstinline|IOException|, but the implementation of that method for \lstinline|StringBuilder| will never throw it and does not even declare it.
\begin{lstlisting}
…
final Appendable appendable = new StringBuilder();
try
{
    appendable.append( "Some Text" );
}
catch( final IOException e )
{
    throw new ImpossibleExceptionError( "Append to StringBuilder", e );
}
…
\end{lstlisting}

Again, a sample for \lstinline|ImpossibleExceptionError| can be found in \autocite{TQUADRAT_ORG_FOUNDATION_IMPOSSIBLEEXCEPTIONERROR}.

Sometimes a method wants to act upon an error condition without handling the exception. In this case, it may catch the exception, does whatever necessary, then throw it again. An example:
\begin{lstlisting}
public final void myMethod() throw IOException
{
    final Predicate<String> isNetworkError = m_NetworkErrorPattern.asMatchPredicate();
    …
    try
    {
        …
    }
    catch( final IOException e )
    {
        final var message = e.getMessage();
        if( nonNull( message ) && isNetworkError.test( message ) ) ++networkErrors;
        throw e;
    }
    …
}   // myMethod()
\end{lstlisting}

It is usually not a good idea to log an exception and to rethrow or to wrap it afterwards, because this might cause the exception to be logged twice (or even more often). More details about logging can be found in chapter \tqvref{sec:Logging}.

Finally, you handle the exception somehow, by assigning a default value\footnote{see the sample about the conversion from string to number, above.}, or by performing an alternative action. In any case, you need to provide a comment in the \lstinline|catch| why you do not further delegate the handling of the exception.

If an exception was not caught in your code, it ”bubbles up” until it reaches the main method for the current thread (that is \lstinline|main()| for the main or program thread, and \lstinline|run()|\autocite{ORACLE_DOC_RUNNABLE:run} for any other thread). If this method will not handle the exception, the thread will be aborted. If this thread is the last non-deamon thread, the JVM will be terminated, too.


%==============================================================================
% End of file!!
%==============================================================================

